\chapter{RTL 구현}\label{ch:rtl}

이 장은 봄베를 디지털 회로로 옮기는 과정을 다룬다. FPGA 구현은 이 프로젝트에서 두 가지 역할을 한다.
첫째, 그 자체가 세 번째 봄베다. 기계식과 릴레이식 옆에 놓여, 움직이는 부품이 하나도 없는 전자 회로가
같은 메뉴에서 같은 정지를 훨씬 빨리 내는 것을 보여 준다. 드럼이 도는 모습이 없으므로 내부 상태(676가닥의
선, 드럼 위치, 정지)를 눈에 보이게 표시하는 장치가 따로 필요하다. 둘째, 기계식과 릴레이식 실물 뒤에서
구동을 제어하고 실물이 낸 정지를 기준 모델의 결과와 대조한다.

회로는 SystemVerilog로 기술했고 두 벌이 있다.
하나는 소프트웨어 모델을 가장 직접적으로 옮긴 \emph{기준 RTL}(\code{rtl/src})로, 시뮬레이션에서만 쓴다.
다른 하나는 실제 FPGA에 올리기 위해 구조를 바꾼 \emph{FPGA 엔진}(\code{rtl/fpga})이다. 기준 RTL과 그 검증은
완성되었으므로 자세히 적고, FPGA 엔진은 현재 조정 중이므로 구조의 개요만 적고 나머지는 자리를 잡아 둔다.

\section{전기 회로를 논리 회로로: 왜 클럭이 필요한가}\label{sec:why-clock}

실제 봄베의 배선망은 전선이다. 전압을 걸면 전류는 이어진 모든 선으로 사실상 즉시 퍼지고, 드럼이 다음
위치로 움직이면 접점이 끊어져 전류는 사라진다. 이것을 논리 게이트로 그대로 옮기면 어떻게 될까.

``케이블 $a$의 선이 통전되면 스크램블러 건너편 케이블 $b$의 선도 통전된다''를 게이트로 쓰면 $b$의 선은
$a$의 선들의 OR이고, 연결이 양방향이므로 $a$의 선은 다시 $b$의 선들의 OR이다. 출력이 게이트 몇 개를 거쳐
자기 입력으로 돌아오는 \term{조합 루프}{combinational loop}가 생긴다(그림~\ref{fig:comb-loop} 왼쪽).
조합 루프는 두 가지 이유로 쓸 수 없다.

\begin{itemize}
\item \textbf{꺼지지 않는다.} OR 게이트의 고리는 한번 1이 되면 입력을 거두어도 서로를 1로 붙들고 있다.
전선의 전류는 접점이 떨어지면 사라지지만, 게이트의 고리는 의도하지 않은 기억 소자가 된다. 다음 드럼
위치를 시험하려면 676개의 선을 모두 0으로 되돌려야 하는데 그럴 방법이 없다.
\item \textbf{설계 도구가 다루지 못한다.} 디지털 회로의 설계 도구는 신호가 레지스터에서 출발해 게이트를
거쳐 다음 레지스터에 도착하는 데 걸리는 시간을 계산해 회로가 제 속도로 동작하는지 확인한다. 고리에는
출발점과 도착점이 없어 이 계산을 할 수 없다.
\end{itemize}

해결은 고리를 \term{레지스터}{register}로 끊는 것이다(그림~\ref{fig:comb-loop} 오른쪽). 676개의 선의
상태를 676비트 레지스터에 담고, 조합 논리는 ``현재 상태에서 \emph{한 걸음} 더 퍼진 상태''만 계산한다.
클럭이 한 번 뛸 때마다 그 결과를 레지스터에 다시 담는다. 상태가 더 변하지 않으면 전류가 다 퍼진 것이다.
새 위치를 시작할 때는 레지스터를 0으로 지우고 전압을 건 선 하나만 1로 둔다.

\begin{figure}[ht]
\centering
\begin{tikzpicture}[x=1cm,y=1cm]
  % 왼쪽: 조합 루프
  \node[font=\small\bfseries] at (2.2,4.0) {그대로 옮긴 경우 (쓸 수 없음)};
  \node[cloud,minimum width=24mm] (fa) at (2.2,2.0) {$a$에서 $b$로\\(OR 망)};
  \node[cloud,minimum width=24mm] (fb) at (2.2,0.2) {$b$에서 $a$로\\(OR 망)};
  \draw[->,cred,line width=0.9pt] (fa.east) -- ++(0.9,0) |- (fb.east);
  \draw[->,cred,line width=0.9pt] (fb.west) -- ++(-0.9,0) |- (fa.west);
  \node[note,text=cred] at (2.2,-0.9) {출력이 입력으로 되돌아온다.\\한번 1이 되면 꺼지지 않는다.};
  % 오른쪽: 레지스터로 끊은 구조
  \begin{scope}[xshift=7.2cm]
  \node[font=\small\bfseries] at (3.0,4.0) {레지스터로 끊은 구조};
  \node[reg,minimum height=18mm,minimum width=14mm] (r) at (0.6,1.1) {\code{live\_q}\\676비트};
  \node[cloud,minimum width=40mm,minimum height=18mm] (f) at (4.3,1.1) {한 걸음 퍼뜨리기\\{\footnotesize 모든 스크램블러 양방향}\\{\footnotesize + 대각선 보드}};
  \draw[->] (r.east) -- (f.west);
  \draw[->] (f.east) -- ++(0.7,0) node[above,pos=0.5,font=\footnotesize\ttfamily] {} |- (4.3,-0.6) node[pos=0.25,right,font=\footnotesize] {\code{live\_n}} -| (r.south);
  \draw (r.south west) ++(0.25,0) -- ++(0.15,0.2) -- ++(0.15,-0.2);
  \node[note] at (0.6,-0.1) {};
  \node[blkn,minimum width=26mm] (cmp) at (4.3,2.9) {\code{live\_n == live\_q} ?};
  \draw[->] (f.north) -- (cmp.south);
  \draw[->] (cmp.east) -- ++(0.8,0) node[right,font=\footnotesize] {끝};
  \node[note] at (3.0,-1.25) {클럭마다 한 걸음. 변화가 없으면 고정점.};
  \end{scope}
\end{tikzpicture}
\caption{전류의 확산을 논리 회로로 옮기는 두 방법. 왼쪽은 조합 루프가 되어 쓸 수 없다. 오른쪽은 상태를
레지스터에 담고 클럭마다 한 걸음씩 퍼뜨린다.}
\label{fig:comb-loop}
\end{figure}

이 구조가 실제 기계, 그리고 소프트웨어 모델과 \emph{같은 답}을 낸다는 것은 \ref{ch:software}장에서 본
성질에서 나온다. 구하는 것은 전압을 건 선과 이어진 선들의 집합이고, 이 집합은 선을 살피는 순서와
무관하다. 소프트웨어는 선을 한 가닥씩 살피고 회로는 모든 선을 동시에 한 걸음씩 살피지만 도달하는 상태는
같다. 회로가 쓰는 클럭 수는 전압을 건 선에서 가장 먼 선까지의 걸음 수로 정해진다.

\section{기준 RTL의 구조}

기준 RTL은 세 모듈이다(그림~\ref{fig:rtl-blocks}). 구성은 실제 기계의 부분과 같다.

\begin{figure}[ht]
\centering
\begin{tikzpicture}[x=1cm,y=1cm]
  \draw[rounded corners=2pt,draw=cgray,fill=clight] (-0.3,-2.9) rectangle (15.0,3.1);
  \node[anchor=north west,font=\small\bfseries] at (-0.2,3.05) {\code{bombe\_sweep}: 드럼 구동과 정지 출력};
  \node[blko,minimum width=26mm,minimum height=12mm] (cnt) at (1.3,1.1) {위치 계수기\\{\footnotesize $g,l,m,r$ (26진)}};
  \node[blko,minimum width=26mm,minimum height=12mm] (fsm) at (1.3,-1.3) {제어\\{\footnotesize IDLE/RUN/EMIT}};
  \draw[rounded corners=2pt,draw=cblue,fill=white] (4.6,-2.4) rectangle (12.6,2.4);
  \node[anchor=north west,font=\small\bfseries] at (4.7,2.35) {\code{bombe\_core}: 한 위치의 판정};
  \node[blkb,minimum width=22mm,minimum height=12mm] (scr) at (6.0,1.1) {\code{scrambler}\\{\footnotesize × 선의 수}};
  \node[reg,minimum width=17mm,minimum height=12mm] (perm) at (8.6,1.1) {\code{perm\_q}\\{\footnotesize 치환표}};
  \node[blkb,minimum width=22mm,minimum height=12mm] (spr) at (11.2,1.1) {전파\\{\footnotesize 조합 논리}};
  \node[reg,minimum width=22mm,minimum height=10mm] (live) at (11.2,-1.3) {\code{live\_q}\\{\footnotesize 676비트}};
  \draw[->] (cnt) -- node[above,font=\scriptsize,text=cgray] {위치} (scr);
  \draw[->] (scr) -- (perm);
  \draw[->] (perm) -- (spr);
  \draw[->] ([xshift=3mm]spr.south) -- ([xshift=3mm]live.north);
  \draw[->] ([xshift=-3mm]live.north) -- ([xshift=-3mm]spr.south);
  \draw[->] ([yshift=2.5mm]fsm.east) -- node[above,font=\scriptsize,text=cgray] {start} ([yshift=2.5mm]4.6,-1.3);
  \draw[<-] ([yshift=-2.5mm]fsm.east) -- node[below,font=\scriptsize,text=cgray] {done, stop} ([yshift=-2.5mm]4.6,-1.3);
  \draw[->] (fsm) -- (cnt);
  \node[blkr,minimum width=15mm,minimum height=12mm] (out) at (13.8,-1.3) {정지\\출력};
  \draw[->] (live) -- (out);
  \node[font=\scriptsize,text=cgray,below=1mm of out] {valid/ready};
\end{tikzpicture}
\caption{기준 RTL의 구성. \code{scrambler}는 조합 논리로 된 치환표, \code{bombe\_core}는 한 드럼 위치에서
전류를 퍼뜨려 정지 여부를 판정하고, \code{bombe\_sweep}은 모든 위치를 차례로 돌며 정지를 내보낸다.}
\label{fig:rtl-blocks}
\end{figure}

\paragraph{\code{scrambler}.} 드럼 한 벌이다. 로터 순서와 코어 위치를 입력으로 받아 26글자의 치환표를
내놓는다. 클럭이 없는 순수한 조합 논리이며, 신호 경로는 소프트웨어의 스크램블러와 같다(오른쪽, 가운데,
왼쪽, 반사판, 다시 역순). 배선표는 소프트웨어와 같은 원본 파일에서 생성한 상수 표다.

\paragraph{\code{bombe\_core}.} 한 드럼 위치의 판정을 맡는다. 메뉴의 선마다 스크램블러를 하나씩 두고,
선의 크립 위치만큼 오른쪽 코어 위치를 더해 준다. 시작 신호가 오면 치환표들을 레지스터에 붙잡고, 676비트
상태를 지운 뒤 전압을 건 선 하나만 1로 둔다. 그다음 클럭마다 한 걸음씩 퍼뜨린다. 한 걸음은 메뉴의 모든
선을 양방향으로, 그리고 대각선 보드를 한꺼번에 적용한 것이다. 상태가 변하지 않거나 테스트 레지스터의
26선이 모두 통전되면 끝낸다. 테스트 레지스터가 26선 모두는 아니면 정지다.

\paragraph{\code{bombe\_sweep}.} 드럼 구동 기구와 정지 릴레이의 역할이다. 26진 자리올림 계수기로 드럼
위치를 하나씩 올리며 \code{bombe\_core}를 구동하고, 정지가 나오면 위치와 테스트 레지스터를 내보낸다.
탐색 순서는 소프트웨어와 같게(오른쪽 로터가 가장 안쪽) 맞췄다. 그래야 정지 목록을 순서까지 포함해
원소 단위로 비교할 수 있다.

\section{한 위치의 동작}

그림~\ref{fig:core-timing}은 \code{bombe\_core}가 한 위치를 판정하는 동안의 신호다. 시작 신호를 본 다음
클럭에서 전압을 건 선 하나만 통전된 상태 $L_0$이 레지스터에 담긴다. 이후 클럭마다 $L_1, L_2, \ldots$로
넓어지다가, 조합 논리가 계산한 다음 상태가 현재 상태와 같아지면(그림에서는 $L_3$) 다음 클럭에 완료
신호를 낸다. 마지막 한 클럭은 ``더 변하지 않는다''는 것을 확인하는 데 쓰인다.

\begin{figure}[ht]
\centering
\begin{tikztimingtable}[timing/slope=0.1,timing/coldist=4pt,xscale=1.9,yscale=1.3,semithick,
    timing/d/background/.style={fill=cbluebg}]
  \code{clk}      & 11{.5L .5H} \\
  \code{start}    & .5L 1H 9.5L \\
  상태            & 1.5D{IDLE} 4D{SPREAD} 5.5D{IDLE} \\
  \code{live\_q}  & 1.5U 1D{$L_0$} 1D{$L_1$} 1D{$L_2$} 6.5D{$L_3$} \\
  \code{live\_n}  & 1.5U 1D{$L_1$} 1D{$L_2$} 7.5D{$L_3$} \\
  \code{done}     & 5.5L 1H 4.5L \\
\extracode
  \begin{pgfonlayer}{background}
    \vertlines[cgray!50,dashed]{0.5,1.5,...,10.5}
  \end{pgfonlayer}
  \node[note,anchor=north] at (5.0,-11.4) {$\uparrow$ \code{live\_n == live\_q}: 고정점};
\end{tikztimingtable}
\caption{\code{bombe\_core}의 동작(개념도). $L_i$는 $i$걸음 퍼진 뒤의 676비트 상태다. 세로 점선은 클럭의
상승 에지다. 실제 걸음 수는 메뉴와 위치에 따라 다르다.}
\label{fig:core-timing}
\end{figure}

시뮬레이션에서 측정한 값은 다음과 같다. 구동 중에는 테스트 레지스터가 다 차는 즉시 끝내므로 위치당 평균
약 12클럭이 든다. 끝까지 퍼뜨리는 경우의 걸음 수는 13회에서 25회 사이였다.

\section{드럼 구동과 정지 출력}

\code{bombe\_sweep}의 제어는 상태가 셋인 상태 기계다(그림~\ref{fig:sweep-fsm}). 정지가 나오면 EMIT 상태에서
위치와 레지스터를 내놓고, 받는 쪽이 받았다는 신호를 줄 때까지 기다린다. 받는 쪽이 준비되지 않으면 탐색이
멈춰 선다. 실제 봄베가 정지한 뒤 조작자가 값을 적고 다시 출발시킬 때까지 서 있던 것과 같다.

\begin{figure}[ht]
\centering
\begin{tikzpicture}[x=1cm,y=1cm,
   st/.style={draw,circle,minimum size=15mm,fill=cbluebg,draw=cblue,font=\small\ttfamily}]
  \node[st] (idle) at (0,0) {IDLE};
  \node[st] (run) at (5,0) {RUN};
  \node[st,fill=credbg,draw=cred] (emit) at (10,0) {EMIT};
  \draw[->] (idle) to[bend left=18] node[above,note] {start: 위치 $\leftarrow 0$, 코어 시작} (run);
  \draw[->] (run) to[bend left=18] node[above,note] {코어 완료, 정지} (emit);
  \draw[->] (emit) to[bend left=18] node[below,note] {받는 쪽이 수락: 다음 위치} (run);
  \draw[->] (run) to[bend left=18] node[below,note] {마지막 위치 완료} (idle);
  \draw[->] (run) to[out=120,in=60,looseness=5] node[above,note] {코어 완료, 통과: 다음 위치} (run);
  \draw[->] (emit) to[out=30,in=-30,looseness=5] node[right,note] {대기} (emit);
\end{tikzpicture}
\caption{\code{bombe\_sweep}의 상태 기계. 마지막 위치가 정지인 경우에는 EMIT에서 바로 IDLE로 간다(그림에서는
생략).}
\label{fig:sweep-fsm}
\end{figure}

\section{검증}\label{sec:rtl-verify}

회로의 검증은 한 문장으로 요약된다. \emph{같은 입력에 대해 Rust 코어와 같은 출력을 내는가.}
이를 위해 Rust 코어가 내보낸 벡터를 Python으로 쓴 테스트벤치(cocotb)가 회로에 넣고, 회로의 출력을 코어의
기대값과 비교한다(그림~\ref{fig:verify}). 시뮬레이터는 Verilator를 주로 쓰고 Icarus Verilog로도 일부를
확인했다. 서로 다른 시뮬레이터에서 같은 결과가 나오는지 보는 것은 회로 기술이 특정 도구의 해석에 기대고
있지 않은지 확인하는 수단이다.

\begin{figure}[ht]
\centering
\begin{tikzpicture}[x=1cm,y=1cm]
  \node[blkb,minimum width=30mm,minimum height=12mm] (core) at (0,1.2) {Rust 코어\\(기준 모델)};
  \node[blko,minimum width=30mm,minimum height=12mm] (vec) at (4.6,1.2) {검증 벡터\\메뉴, 위치, 기대 상태};
  \node[blkg,minimum width=30mm,minimum height=12mm] (tb) at (9.2,1.2) {테스트벤치\\(cocotb, Python)};
  \node[blkn,minimum width=30mm,minimum height=12mm] (dut) at (9.2,-1.0) {회로\\(시뮬레이터)};
  \node[blkr,minimum width=30mm,minimum height=12mm] (cmp) at (13.4,0.1) {비교\\676비트, 정지 목록};
  \draw[->] (core) -- (vec);
  \draw[->] (vec) -- (tb);
  \draw[->] ([xshift=-4mm]tb.south) -- node[left,note] {입력} ([xshift=-4mm]dut.north);
  \draw[->] ([xshift=4mm]dut.north) -- node[right,note] {출력} ([xshift=4mm]tb.south);
  \draw[->] (tb.east) -- ++(0.5,0) |- ([yshift=2mm]cmp.west);
  \draw[->,dashed] (vec.south) -- ++(0,-2.6) -| node[pos=0.12,above,note] {기대값} (cmp.south);
\end{tikzpicture}
\caption{검증의 흐름. 회로는 코어가 내보낸 기대값과 비트 단위로 비교된다.}
\label{fig:verify}
\end{figure}

시험은 모듈의 층마다 따로 두었다(표~\ref{tab:rtl-tests}). 아래층이 맞다는 것을 확인한 뒤 위층을 시험해야
실패했을 때 원인을 좁힐 수 있기 때문이다.

\begin{table}[ht]
\centering
\caption{기준 RTL의 시험. Verilator에서 일곱 항목이 모두 통과한다(전체 89초).}
\label{tab:rtl-tests}
\small
\begin{tabularx}{\linewidth}{l>{\raggedright\arraybackslash}X>{\raggedright\arraybackslash}X}
\toprule
대상 & 내용 & 비교하는 것 \\
\midrule
\code{scrambler} & 무작위 로터 순서와 위치 3,000건(모든 반사판, 그리스 로터 유무), 오른쪽 로터 전수,
  \texttt{AAAAA}$\to$\texttt{BDZGO} & 26글자 치환표 \\
\code{bombe\_core} & 벡터 파일 6개(넘어감 가설, 대각선 보드 없음, M4 포함) & 676선 상태 전체, 정지 여부, 테스트 레지스터 \\
\code{bombe\_sweep} & Enigma~I 전체 스윕 4종, M4 스윕의 앞 40개 정지 & 정지 목록(순서 포함), 시험한 위치 수 \\
\bottomrule
\end{tabularx}
\end{table}

두 가지 세부 사항이 검증의 강도를 높인다. 첫째, 테스트벤치는 쓰지 않는 입력(비활성화된 메뉴 선의 칸,
3로터에서의 그리스 로터 번호)에 무작위 값을 넣는다. 회로가 쓰지 않아야 할 입력에 실수로 의존하고 있으면
이때 드러난다. 둘째, 스윕 시험은 정지 출력의 수락 신호를 무작위로 늦춘다. 받는 쪽이 늦을 때 정지가
사라지거나 두 번 나오지 않는지, 기다리는 동안 출력값이 유지되는지를 본다.

\section{FPGA 엔진의 개요}\label{sec:fpga-overview}

\subsection{기준 RTL을 그대로 올릴 수 없는 이유}

기준 RTL은 메뉴의 선마다 스크램블러를 하나씩 두고, 드럼 위치가 바뀔 때마다 모든 스크램블러가 26글자의
치환을 한꺼번에 다시 계산한다. 시뮬레이션에서는 문제가 없지만 실제 회로로 만들면 이 계산 하나하나가
논리 소자를 차지한다. 합성해 보면 스크램블러 하나가 LUT(FPGA의 기본 논리 소자) 19,717개를 쓴다.
선이 16개인 메뉴라면 스크램블러만으로 약 31만 5천 개가 되어, 대상 FPGA(XC7A200T, LUT 약 13만 4천 개)에
들어가지 않는다.

\subsection{바뀌는 것만 다시 계산한다}

해결의 실마리는 \ref{sec:sw-tables}절과 같다. 가운데와 왼쪽 로터가 고정된 동안에는 오른쪽 로터의
26가지 위치에 대한 치환표 26개만 있으면 되고, 그 표는 가운데 로터가 바뀔 때에만 다시 만들면 된다.
실제 기계의 느린 드럼이 빠른 드럼 한 바퀴에 한 번만 움직이는 것과 같은 구조다. FPGA 엔진은 이 표를
여러 클럭에 걸쳐 조금씩 계산해 FPGA 안의 작은 메모리에 넣어 두고, 위치마다 표에서 읽어 쓴다. 논리 소자를
시간과 맞바꾼 것이다(그림~\ref{fig:fpga-blocks}).

\begin{figure}[ht]
\centering
\begin{tikzpicture}[x=1cm,y=1cm]
  \node[blkg,minimum width=22mm,minimum height=14mm] (host) at (-0.2,-0.1) {호스트\\{\scriptsize PC 또는 보드의 ARM}};
  \node[blkn,minimum width=22mm,minimum height=14mm] (br) at (2.8,-0.1) {브리지\\{\scriptsize UART / AXI4-Lite}};
  \draw[rounded corners=2pt,draw=cgray,fill=clight] (5.9,-3.1) rectangle (14.5,2.0);
  \node[anchor=north west,font=\small\bfseries] at (6.0,1.95) {\code{bombe\_array}: 엔진 $N$대};
  \draw[rounded corners=2pt,draw=cblue,fill=white] (6.2,-2.7) rectangle (14.2,1.2);
  \node[anchor=north west,font=\small\bfseries] at (6.3,1.15) {\code{bombe\_engine} (로터 순서 하나를 스윕)};
  \node[blko,minimum width=21mm] (cfg) at (7.5,-0.1) {설정 레지스터\\{\scriptsize 메뉴, 로터 순서}};
  \node[blkb,minimum width=21mm] (tbl) at (10.2,-0.1) {치환표 계산\\{\scriptsize 파이프라인}};
  \node[reg,minimum width=20mm] (ram) at (12.8,-0.1) {표 메모리\\{\scriptsize 분산 RAM}};
  \node[blkb,minimum width=21mm] (spr) at (10.2,-1.8) {전파\\{\scriptsize 676비트 상태}};
  \node[blkr,minimum width=21mm] (fifo) at (7.5,-1.8) {정지 FIFO};
  \draw[->] (cfg) -- (tbl); \draw[->] (tbl) -- (ram);
  \draw[->] (ram.south) |- (spr.east);
  \draw[->] (spr) -- (fifo);
  \draw[<->] (host) -- (br);
  \draw[<->] (br) -- node[above,font=\scriptsize,text=cgray] {레지스터 버스} (6.2,-0.1);
\end{tikzpicture}
\caption{FPGA 엔진의 블록 수준 구성. 엔진 하나가 로터 순서 하나를 맡고, 여러 엔진이 레지스터 버스 뒤에
나란히 놓인다. 호스트는 메뉴와 로터 순서를 써넣고 정지를 읽어 간다.}
\label{fig:fpga-blocks}
\end{figure}

엔진의 정지 목록은 기준 RTL과 Rust 코어의 정지 목록과 순서까지 같아야 한다. 시뮬레이션에서 엔진 두 대를
동시에 구동한 Enigma~I 전체 스윕의 정지 목록이 Rust 코어와 순서까지 일치하는 것을 확인했다.

\subsection{회로의 크기는 계산의 양이 아니라 기술한 구조가 정한다}

FPGA 엔진의 첫 구현은 합성 결과 LUT 약 40만 4천 개, 대상 FPGA 용량의 300\%였다. 원인은 계산의 양이 아니라
676비트 상태에 접근하는 방식이었다.

\begin{lesson}{676비트를 정수 번호로 고르면 거대한 선택 회로가 된다}
전파 논리는 ``케이블 $a$의 $y$번 선이 통전되어 있으면 케이블 $b$의 $x$번 선을 통전시킨다''를 메뉴의 선마다,
글자마다 한다. 이를 676비트 벡터의 ``$26a + y$번째 비트''로 적으면 소프트웨어에서는 배열 접근 한 번이다.
회로에서는 다르다. $a$와 $y$가 실행 중에 정해지는 값이므로, 이 한 줄은 \emph{676개의 비트 가운데 하나를
골라 내는 선택 회로}가 된다. 메뉴의 선 16개, 글자 26개, 양방향이면 이런 선택 회로가 832개 필요하고,
결과를 써넣는 쪽에도 같은 수의 ``676곳 중 한 곳에 쓰기'' 회로가 생긴다.

구조를 케이블 단위로 바꾸자 같은 기능이 LUT 41,040개(용량의 30.5\%)로 줄었다. 메뉴의 선마다 먼저
케이블 두 개를 통째로 고르고(26개 가운데 하나를 고르는 26비트 폭의 선택), 그 안에서 치환에 따라 선을
옮긴 뒤(선마다 26개 가운데 하나를 고르는 선택), 결과를 목적지 케이블에 OR로 모은다. 목적지 케이블을
가리키는 신호는 구동을 시작할 때 한 번만 만들어 둔다. 676개 중 하나를 고르는 일을 26개 중 하나를 고르는
일 두 번으로 나눈 것이다.
\end{lesson}

\begin{figure}[ht]
\centering
\begin{tikzpicture}
\begin{axis}[
    xbar, width=0.78\linewidth, height=50mm,
    xmin=0, xmax=480000, bar width=7mm, bar shift=0pt,
    ymin=-0.7, ymax=1.9, clip=false,
    ytick={0,1}, yticklabels={케이블 단위 구조,첫 구현},
    y tick label style={font=\small},
    xlabel={LUT 수 (Vivado 합성, XC7A200T)}, xlabel style={font=\small},
    scaled x ticks=false,
    xtick={0,100000,200000,300000,400000},
    xticklabels={0,10만,20만,30만,40만},
    x tick label style={font=\small},
    axis lines*=left,
  ]
  \addplot[fill=cblue!70,draw=cblue] coordinates {(41040,0)};
  \addplot[fill=cred!70,draw=cred] coordinates {(404000,1)};
  \node[anchor=west,font=\small] at (axis cs:41040,0) {\,41,040 (30.5\%)};
  \node[anchor=west,font=\small] at (axis cs:404000,1) {\,약 40만 4천 (300\%)};
  \draw[dashed,line width=0.9pt,cgray] (axis cs:134000,-0.7) -- (axis cs:134000,1.75)
     node[above,font=\footnotesize,text=cgray] {FPGA 용량 약 13만 4천};
\end{axis}
\end{tikzpicture}
\caption{같은 기능을 두 가지 구조로 기술했을 때의 회로 크기.}
\label{fig:lut-bars}
\end{figure}

케이블 단위 구조에서 나머지 자원은 LUT RAM 2,858개, 플립플롭 7,338개다. 치환표를 계산하는 안쪽
파이프라인은 처음에 7단이었으나 단 사이의 논리가 너무 길어 목표 클럭을 맞추지 못했고, 14단으로 나누어
해결했다. 이 수치들은 2026년 11월 14일 기준이며 엔진은 타이밍을 맞추는 조정이 진행 중이다.

\begin{todobox}{FPGA 엔진의 상세}
엔진의 설계가 확정된 뒤 다음 내용을 쓴다.
\begin{itemize}
\item 배선표 적재, 안쪽 치환 파이프라인, 오른쪽 로터 26위치 치환표(병렬 레인)의 구조와 타이밍 그림
\item 메뉴의 선마다 분산 RAM을 두는 이유와 읽기 경로
\item 케이블 단위 전파 회로의 구조 그림(케이블 선택, 치환, 원핫 OR)
\item 고정점 판정과 조기 종료의 구현, 위치당 클럭 수의 분포
\item 정지 FIFO와 역압(backpressure), 레지스터 맵
\item 확정된 자원 사용량과 타이밍 결과(배치·배선 후), 엔진 수에 따른 변화
\end{itemize}
\end{todobox}

\begin{todobox}{보드 통합}
\begin{itemize}
\item Nexys Video(Artix-7 XC7A200T): UART 브리지, 최상위 모듈, 제약 파일, 빌드 절차
\item PYNQ(Zynq XC7Z020): AXI4-Lite 브리지와 블록 디자인. 현재 설치된 도구에 이 디바이스 지원이 없어
빌드하지 못한 상태다
\item 호스트 소프트웨어: 여러 엔진에 로터 순서를 분배하고 결과를 Rust 코어와 대조하는 절차
\end{itemize}
\end{todobox}

\begin{todobox}{FPGA 결과}
보드에서의 실측 결과(클럭 주파수, 엔진 수, 로터 순서 하나의 스윕 시간, 전체 탐색 시간, 소프트웨어 및
실제 봄베와의 비교)는 \ref{ch:results}장에 함께 정리한다. 전시물로서의 구성(내부 상태를 LED나 화면으로
보여 주는 시각화 장치, 기계식·릴레이식과의 속도 비교 표시)도 설계가 정해지면 이 장에 적는다.
\end{todobox}
